% =============================================================================
%  CHAPTER 2: Text Formatting & Fonts
% =============================================================================

\chapter{Text Formatting \& Fonts}

\section{Font Styles}

\latex{} provides command pairs for font switching — one with \emph{declaration syntax} (affects all following text) and one with \emph{command syntax} (takes an argument).

\begin{table}[H]
\centering
\begin{tabular}{lll}
\toprule
\textbf{Style} & \textbf{Declaration} & \textbf{Command} \\
\midrule
\textbf{Bold}          & \verb|\bfseries|  & \verb|\textbf{...}| \\
\textit{Italic}        & \verb|\itshape|   & \verb|\textit{...}| \\
\textsl{Slanted}       & \verb|\slshape|   & \verb|\textsl{...}| \\
\texttt{Monospace}     & \verb|\ttfamily|  & \verb|\texttt{...}| \\
\textsf{Sans-serif}    & \verb|\sffamily|  & \verb|\textsf{...}| \\
\textsc{Small Caps}    & \verb|\scshape|   & \verb|\textsc{...}| \\
\textmd{Medium}        & \verb|\mdseries|  & \verb|\textmd{...}| \\
Upright                & \verb|\upshape|   & \verb|\textup{...}| \\
\bottomrule
\end{tabular}
\caption{Font style commands}
\end{table}

\subsection{Examples}

\begin{itemize}
    \item \textbf{Bold text} with \verb|\textbf{Bold text}|
    \item \textit{Italic text} with \verb|\textit{Italic text}|
    \item \texttt{Typewriter text} with \verb|\texttt{Typewriter text}|
    \item \textsf{Sans-serif text} with \verb|\textsf{Sans-serif text}|
    \item \textsc{Small Capitals} with \verb|\textsc{Small Capitals}|
    \item \textsl{Slanted text} with \verb|\textsl{Slanted text}|
    \item \emph{Emphasized text} — \verb|\emph{}| toggles italic inside roman and roman inside italic
    \item \emph{This is \emph{nested} emphasis} — \verb|\emph{This is \emph{nested} emphasis}|
\end{itemize}

% ------------------------------------------------------------------
\section{Font Sizes}

Size commands are \emph{declarations} — they affect everything until the end of the group:

\begin{table}[H]
\centering
\begin{tabular}{ll}
\toprule
\textbf{Command} & \textbf{Result} \\
\midrule
\verb|\tiny|         & {\tiny Sample text} \\
\verb|\scriptsize|   & {\scriptsize Sample text} \\
\verb|\footnotesize| & {\footnotesize Sample text} \\
\verb|\small|        & {\small Sample text} \\
\verb|\normalsize|   & {\normalsize Sample text} \\
\verb|\large|        & {\large Sample text} \\
\verb|\Large|        & {\Large Sample text} \\
\verb|\LARGE|        & {\LARGE Sample text} \\
\verb|\huge|         & {\huge Sample text} \\
\verb|\Huge|         & {\Huge Sample text} \\
\bottomrule
\end{tabular}
\caption{Font size commands (relative to base document font size)}
\end{table}

\begin{lstlisting}[language={[LaTeX]TeX}, caption=Font size usage]
{\large This text is large.}
{\Large\bfseries This text is large and bold.}
\end{lstlisting}

% ------------------------------------------------------------------
\section{Text Alignment}

\subsection{Centering}

\begin{lstlisting}[language={[LaTeX]TeX}, caption=Centering text]
% Environment form:
\begin{center}
    This text is centered.
\end{center}

% Declaration form (affects rest of group):
{\centering This text is centered.\par}
\end{lstlisting}

\begin{center}
    This is centered text using the \texttt{center} environment.
\end{center}

\subsection{Flush Left (Ragged Right)}

\begin{lstlisting}[language={[LaTeX]TeX}, caption=Left-aligned text]
\begin{flushleft}
    This text is left-aligned.
\end{flushleft}
\end{lstlisting}

\subsection{Flush Right (Ragged Left)}

\begin{lstlisting}[language={[LaTeX]TeX}, caption=Right-aligned text]
\begin{flushright}
    This text is right-aligned.
\end{flushright}
\end{lstlisting}

\begin{flushright}
    This is right-aligned text.
\end{flushright}

% ------------------------------------------------------------------
\section{Line \& Paragraph Spacing}

\begin{lstlisting}[language={[LaTeX]TeX}, caption=Line spacing]
% Single spacing (default)
% Double spacing:
\usepackage{setspace}
\doublespacing

% 1.5 spacing:
\onehalfspacing

% Custom spacing:
\setstretch{1.25}

% Line break (new line, same paragraph):
Line one.\\
Line two.

% Line break with extra space:
Line one.\\[12pt]
Line two with 12pt gap above.

% New paragraph (blank line in source):
Paragraph one.

Paragraph two.

% Horizontal spacing:
Word\hspace{1cm}Word          % 1cm horizontal space
Word\hspace{\fill}Word        % Stretchy space (pushes to margins)
Word\hfill Word               % Shorthand for \hspace{\fill}
\end{lstlisting}

\subsection{Vertical Spacing}

\begin{lstlisting}[language={[LaTeX]TeX}, caption=Vertical spacing]
\vspace{1cm}          % 1cm vertical space
\vspace*{1cm}         % Forced space (even at page top)
\vspace{\fill}        % Stretchy vertical space
\vfill                % Shorthand for \vspace{\fill}
\bigskip              % Large vertical skip (~12pt)
\medskip              % Medium vertical skip (~6pt)
\smallskip            % Small vertical skip (~3pt)
\end{lstlisting}

\subsection{Horizontal Spacing}

\begin{table}[H]
\centering
\begin{tabular}{lll}
\toprule
\textbf{Command} & \textbf{Width} & \textbf{Description} \\
\midrule
\verb|\,|    or \verb|\thinspace|   & $\sim$3/18 quad & Thin space \\
\verb|\:|    or \verb|\medspace|    & $\sim$4/18 quad & Medium space \\
\verb|\;|    or \verb|\thickspace|  & $\sim$5/18 quad & Thick space \\
\verb|\!|                          & $-$3/18 quad    & Negative thin space \\
\verb|\enspace|                     & 0.5em           & En-space \\
\verb|\quad|                        & 1em             & Quad space \\
\verb|\qquad|                       & 2em             & Double quad \\
\verb|\~{}|                         & variable        & Non-breaking space \\
\verb|\hspace{length}|              & variable        & Custom horizontal space \\
\bottomrule
\end{tabular}
\caption{Horizontal spacing commands}
\end{table}

% ------------------------------------------------------------------
\section{Text Decoration \& Color}

\subsection{Underline and Overline}

\begin{lstlisting}[language={[LaTeX]TeX}, caption=Underline and overline]
\underline{Underlined text}
\overline{Overlined text}        % Works in math mode too
\end{lstlisting}

\underline{Underlined text example.}

\subsection{Strikethrough (requires ulem package)}

\begin{lstlisting}[language={[LaTeX]TeX}, caption=Strikethrough]
\usepackage[normalem]{ulem}

\sout{Strikethrough text}
\uwave{Wavy underline}
\xout{Crossed out with /}
\end{lstlisting}

\subsection{Text Color (xcolor package)}

\begin{lstlisting}[language={[LaTeX]TeX}, caption=Using colors]
\textcolor{red}{Red text}
\textcolor{blue}{Blue text}
\textcolor{green!50!black}{Dark green text}

% Color as declaration:
{\color{red} This entire block is red.}

% Define custom colors:
\definecolor{mycolor}{RGB}{102, 178, 255}
\textcolor{mycolor}{Custom colored text}

% Background color (requires xcolor with table option):
\colorbox{yellow}{Highlighted text}
\fcolorbox{red}{yellow}{Box with red border and yellow background}
\end{lstlisting}

Examples: \textcolor{red}{Red}, \textcolor{blue}{Blue}, \textcolor{green!50!black}{Dark Green}, \colorbox{yellow}{Highlighted}

% ------------------------------------------------------------------
\section{Accents and Special Characters}

\begin{table}[H]
\centering
\begin{tabular}{llll}
\toprule
\textbf{Accent} & \textbf{Command} & \textbf{Result} & \textbf{Alt (UTF-8)} \\
\midrule
Grave        & \verb|\`{o}| & \`{o} & \`o \\
Acute        & \verb|\'{o}| & \'{o} & \'o \\
Circumflex   & \verb|\^{o}| & \^{o} & \^o \\
Tilde        & \verb|\~{o}| & \~{o} & \~o \\
Umlaut       & \verb|\"{o}| & \"{o} & \"o \\
Cedilla      & \verb|\c{c}| & \c{c} & \c{c} \\
Ogonek       & \verb|\k{a}| & \k{a} & \\
Macron       & \verb|\={o}| & \={o} & \\
Breve        & \verb|\u{o}| & \u{o} & \\
Dot          & \verb|\.{o}| & \.{o} & \\
Ring         & \verb|\r{a}| & \r{a} & \\
Hacek        & \verb|\v{o}| & \v{o} & \\
Double acute & \verb|\H{o}| & \H{o} & \\
\bottomrule
\end{tabular}
\caption{Accent commands}
\end{table}

% ------------------------------------------------------------------
\section{Dashes and Quotes}

\begin{table}[H]
\centering
\begin{tabular}{llll}
\toprule
\textbf{Type} & \textbf{Source} & \textbf{Result} & \textbf{Usage} \\
\midrule
Hyphen    & \verb|-|   & mother-in-law & Compound words \\
En-dash   & \verb|--|  & pages 1--10   & Ranges \\
Em-dash   & \verb|---| & Yes---indeed   & Interruptions \\
Minus     & \verb|$-$| & $-$           & Mathematical minus \\
\bottomrule
\end{tabular}
\caption{Types of dashes}
\end{table}

\begin{table}[H]
\centering
\begin{tabular}{lll}
\toprule
\textbf{Type} & \textbf{Source} & \textbf{Result} \\
\midrule
Opening double quotes & \verb|``| & ``Hello'' \\
Closing double quotes & \verb|''| & ``Hello'' \\
Opening single quote  & \verb|`|  & `Hello' \\
Closing single quote  & \verb|'|  & `Hello' \\
\bottomrule
\end{tabular}
\caption{Quotation marks}
\end{table}

\note{\latex{} does not use regular keyboard quotes (\texttt{"}). Always use backticks for opening quotes.}
